\documentclass[10pt,a4paper]{amsart}
\usepackage{amsmath,amssymb,amsthm,graphicx}
% Retained preamble for the Stacks Project chapter snapshot in this directory.
% The chapter's own preamble.tex is a copy of the shared Stacks Project preamble; the block
% below is sliced MECHANICALLY, line for line, from it. Nothing was added, removed or
% redefined; the wrapper class, the xr-hyper/hyperref wiring and the externaldocument list
% are omitted because the assembler supplies the document class.
%
% lines taken from preamble.tex: 3, 5, 9, 11, 12, 156, 157, 158, ... (51 lines total)
%
% package coverage restored from the chapter's own preamble (2026-10-13): xy, xy-2cell

\usepackage[all]{xy}
\xyoption{2cell}
\UseAllTwocells

\newenvironment{reference}{\comment}{\endcomment}
\newenvironment{slogan}{\comment}{\endcomment}
\newenvironment{history}{\comment}{\endcomment}
\theoremstyle{plain}
\newtheorem{theorem}[subsection]{Theorem}
\newtheorem{proposition}[subsection]{Proposition}
\newtheorem{lemma}[subsection]{Lemma}
\theoremstyle{definition}
\newtheorem{definition}[subsection]{Definition}
\newtheorem{example}[subsection]{Example}
\newtheorem{exercise}[subsection]{Exercise}
\newtheorem{situation}[subsection]{Situation}
\theoremstyle{remark}
\newtheorem{remark}[subsection]{Remark}
\newtheorem{remarks}[subsection]{Remarks}
\numberwithin{equation}{subsection}
\def\lim{\mathop{\mathrm{lim}}\nolimits}
\def\colim{\mathop{\mathrm{colim}}\nolimits}
\def\Spec{\mathop{\mathrm{Spec}}}
\def\Hom{\mathop{\mathrm{Hom}}\nolimits}
\def\Ext{\mathop{\mathrm{Ext}}\nolimits}
\def\SheafHom{\mathop{\mathcal{H}\!\mathit{om}}\nolimits}
\def\SheafExt{\mathop{\mathcal{E}\!\mathit{xt}}\nolimits}
\def\Sch{\mathit{Sch}}
\def\Mor{\mathop{\mathrm{Mor}}\nolimits}
\def\Ob{\mathop{\mathrm{Ob}}\nolimits}
\def\Sh{\mathop{\mathit{Sh}}\nolimits}
\def\NL{\mathop{N\!L}\nolimits}
\def\CH{\mathop{\mathrm{CH}}\nolimits}
\def\proetale{{pro\text{-}\acute{e}tale}}
\def\etale{{\acute{e}tale}}
\def\QCoh{\mathit{QCoh}}
\def\Ker{\mathop{\mathrm{Ker}}}
\def\Im{\mathop{\mathrm{Im}}}
\def\Coker{\mathop{\mathrm{Coker}}}
\def\Coim{\mathop{\mathrm{Coim}}}
\DeclareMathSymbol{\boxtimes}{\mathbin}{AMSa}{"02}
\def\QCohstack{\mathcal{QC}\!\mathit{oh}}
\def\Cohstack{\mathcal{C}\!\mathit{oh}}
\def\Spacesstack{\mathcal{S}\!\mathit{paces}}
\def\Quotfunctor{\mathrm{Quot}}
\def\Hilbfunctor{\mathrm{Hilb}}
\def\Curvesstack{\mathcal{C}\!\mathit{urves}}
\def\Polarizedstack{\mathcal{P}\!\mathit{olarized}}
\def\Complexesstack{\mathcal{C}\!\mathit{omplexes}}
\def\Pic{\mathop{\mathrm{Pic}}\nolimits}
\def\Picardstack{\mathcal{P}\!\mathit{ic}}
\def\Picardfunctor{\mathrm{Pic}}
\def\Deformationcategory{\mathcal{D}\!\mathit{ef}}
\usepackage{cleveref}
\crefname{originallemma}{original lemma}{original lemmas}
\Crefname{originallemma}{Original lemma}{Original lemmas}
\makeatletter
\newlabel{defos-lemma-choices}{{08S7}{source}{}{sourcepointer-1}{}}
\newlabel{defos-lemma-choices@cref}{{[originallemma][0][]08S7}{[1][0][]source}{}{}{}}
\makeatother
\crefname{originallemma}{original lemma}{original lemmas}
\Crefname{originallemma}{Original lemma}{Original lemmas}
\makeatletter
\newlabel{defos-lemma-huge-diagram}{{08S5}{source}{}{sourcepointer-2}{}}
\newlabel{defos-lemma-huge-diagram@cref}{{[originallemma][0][]08S5}{[1][0][]source}{}{}{}}
\makeatother
\crefname{originallemma}{original lemma}{original lemmas}
\Crefname{originallemma}{Original lemma}{Original lemmas}
\makeatletter
\newlabel{defos-lemma-parametrize-solutions}{{0GPX}{source}{}{sourcepointer-3}{}}
\newlabel{defos-lemma-parametrize-solutions@cref}{{[originallemma][0][]0GPX}{[1][0][]source}{}{}{}}
\makeatother
\crefname{originalremark}{original remark}{original remarks}
\Crefname{originalremark}{Original remark}{Original remarks}
\makeatletter
\newlabel{defos-remark-parametrize-solutions}{{0GPY}{source}{}{sourcepointer-4}{}}
\newlabel{defos-remark-parametrize-solutions@cref}{{[originalremark][0][]0GPY}{[1][0][]source}{}{}{}}
\makeatother
\crefname{originalequation}{original equation}{original equations}
\Crefname{originalequation}{Original equation}{Original equations}
\makeatletter
\newlabel{equation-to-solve}{{08SN}{source}{}{sourcepointer-5}{}}
\newlabel{equation-to-solve@cref}{{[originalequation][0][]08SN}{[1][0][]source}{}{}{}}
\makeatother
\crefname{originalequation}{original equation}{original equations}
\Crefname{originalequation}{Original equation}{Original equations}
\makeatletter
\newlabel{equation-triangle}{{08QS}{source}{}{sourcepointer-6}{}}
\newlabel{equation-triangle@cref}{{[originalequation][0][]08QS}{[1][0][]source}{}{}{}}
\makeatother
\crefname{originallemma}{original lemma}{original lemmas}
\Crefname{originallemma}{Original lemma}{Original lemmas}
\makeatletter
\newlabel{lemma-identify-H0}{{08QF}{source}{}{sourcepointer-7}{}}
\newlabel{lemma-identify-H0@cref}{{[originallemma][0][]08QF}{[1][0][]source}{}{}{}}
\makeatother
\crefname{originallemma}{original lemma}{original lemmas}
\Crefname{originallemma}{Original lemma}{Original lemmas}
\makeatletter
\newlabel{lemma-relation-with-naive-cotangent-complex}{{08RB}{source}{}{sourcepointer-8}{}}
\newlabel{lemma-relation-with-naive-cotangent-complex@cref}{{[originallemma][0][]08RB}{[1][0][]source}{}{}{}}
\makeatother
\title{The obstruction class for lifting an algebra along a surjection and the parametrization of solutions}
\author{The Stacks Project Authors}
\date{stacks-project:cotangent.tex:local-snapshot-20261012; GNU FDL 1.2 or later}
\begin{document}
\maketitle
\setcounter{subsection}{1}
\noindent\textit{Editorial note.} All delivered mathematical text is the author's own, copied from the pinned Stacks Project chapter file cotangent.tex; no statement or proof step was rewritten, reordered or removed.\par\medskip
\noindent\textit{Editorial note.} The statement refers to the situation and the equation-to-solve fixed earlier in the chapter; the proof opens by recalling which identifications it uses and then treats the three parts.\par\medskip
\noindent\textit{Editorial note.} The eight references of the counted spans are recorded as non-load-bearing navigation pointers; the argument that consumes them stays inside the counted proof span.\par\medskip
\noindent\textit{Editorial note.} Preamble: the retained preamble of this chapter (sources/stacks-cotangent/retained-preamble.tex, lines 1-63), containing the Xy-pic setup, the comment environments, the theorem declarations and the macros including \textbackslash{}NL.\par\medskip
\noindent\textit{Allowed standard prerequisites (editorial).}\begin{itemize}
\item The identification of the naive cotangent complex with the truncation of the cotangent complex. Lemma lemma-relation-with-naive-cotangent-complex, used at proof lines 3185-3187. Transfers the deformation-theory statements to the cotangent complex.
\item The identification H\textasciicircum{}0(L\_\{B/A\}) = Omega\_\{B/A\}. Lemma lemma-identify-H0 (583-586), used at proof lines 3186-3188. Makes the deformation-theory identifications apply.
\item The obstruction class and parametrization of solutions in deformation theory. Deformation Theory, Lemma defos-lemma-parametrize-solutions, Remark defos-remark-parametrize-solutions, Lemmas defos-lemma-huge-diagram and defos-lemma-choices, used at proof lines 3188-3199 and 3206-3210. Supply parts (2) and (3) and the class xi' for part (1).
\item The distinguished triangle for A' -> A -> B and its long exact Ext sequence. equation-triangle, used at proof lines 3216-3223. Produces the map whose image is the obstruction.
\end{itemize}\medskip
\ifdefined\hypertarget\hypertarget{sourcepointer-1}{}\fi
\noindent\textit{Source reference.} Original lemma 08S7: focused-seeds/AGP/sources/stacks-snapshot-20261011/tags.txt, source lines 11120--11120. Navigation pointer into the Stacks Project; the standard fact it supplies is applied inside the counted proof span and the argument resting on it stays there.\par\medskip
\ifdefined\hypertarget\hypertarget{sourcepointer-2}{}\fi
\noindent\textit{Source reference.} Original lemma 08S5: focused-seeds/AGP/sources/stacks-snapshot-20261011/tags.txt, source lines 11118--11118. Navigation pointer into the Stacks Project; the standard fact it supplies is applied inside the counted proof span and the argument resting on it stays there.\par\medskip
\ifdefined\hypertarget\hypertarget{sourcepointer-3}{}\fi
\noindent\textit{Source reference.} Original lemma 0GPX: focused-seeds/AGP/sources/stacks-snapshot-20261011/tags.txt, source lines 20842--20842. Navigation pointer into the Stacks Project; the standard fact it supplies is applied inside the counted proof span and the argument resting on it stays there.\par\medskip
\ifdefined\hypertarget\hypertarget{sourcepointer-4}{}\fi
\noindent\textit{Source reference.} Original remark 0GPY: focused-seeds/AGP/sources/stacks-snapshot-20261011/tags.txt, source lines 20843--20843. Navigation pointer into the Stacks Project; the standard fact it supplies is applied inside the counted proof span and the argument resting on it stays there.\par\medskip
\ifdefined\hypertarget\hypertarget{sourcepointer-5}{}\fi
\noindent\textit{Source reference.} Original equation 08SN: focused-seeds/AGP/sources/stacks-snapshot-20261011/tags.txt, source lines 11136--11136. Navigation pointer into the Stacks Project; the standard fact it supplies is applied inside the counted proof span and the argument resting on it stays there.\par\medskip
\ifdefined\hypertarget\hypertarget{sourcepointer-6}{}\fi
\noindent\textit{Source reference.} Original equation 08QS: focused-seeds/AGP/sources/stacks-snapshot-20261011/tags.txt, source lines 11070--11070. Navigation pointer into the Stacks Project; the standard fact it supplies is applied inside the counted proof span and the argument resting on it stays there.\par\medskip
\ifdefined\hypertarget\hypertarget{sourcepointer-7}{}\fi
\noindent\textit{Source reference.} Original lemma 08QF: focused-seeds/AGP/sources/stacks-snapshot-20261011/tags.txt, source lines 11058--11058. Navigation pointer into the Stacks Project; the standard fact it supplies is applied inside the counted proof span and the argument resting on it stays there.\par\medskip
\ifdefined\hypertarget\hypertarget{sourcepointer-8}{}\fi
\noindent\textit{Source reference.} Original lemma 08RB: focused-seeds/AGP/sources/stacks-snapshot-20261011/tags.txt, source lines 11089--11089. Navigation pointer into the Stacks Project; the standard fact it supplies is applied inside the counted proof span and the argument resting on it stays there.\par\medskip
\begin{lemma}
\label{lemma-find-obstruction}
In the situation above we have
\begin{enumerate}
\item There is a canonical element $\xi \in \Ext^2_B(L_{B/A}, N)$
whose vanishing is a sufficient and necessary condition for the existence
of a solution to (\ref{equation-to-solve}).
\item If there exists a solution, then the set of
isomorphism classes of solutions is principal homogeneous under
$\Ext^1_B(L_{B/A}, N)$.
\item Given a solution $B'$, the set of automorphisms of $B'$
fitting into (\ref{equation-to-solve}) is canonically isomorphic
to $\Ext^0_B(L_{B/A}, N)$.
\end{enumerate}
\end{lemma}

\begin{proof}
Via the identifications $\NL_{B/A} = \tau_{\geq -1}L_{B/A}$
(Lemma \ref{lemma-relation-with-naive-cotangent-complex}) and
$H^0(L_{B/A}) = \Omega_{B/A}$ (Lemma \ref{lemma-identify-H0})
we have seen parts (2) and (3) in
Deformation Theory, Lemmas \ref{defos-lemma-huge-diagram} and
\ref{defos-lemma-choices}.

\medskip\noindent
Proof of (1). Roughly speaking, this follows from the discussion in 
Deformation Theory, Remark \ref{defos-remark-parametrize-solutions}
by replacing the naive cotangent complex by the full cotangent complex.
Here is a more detailed explanation. By
Deformation Theory, Lemma \ref{defos-lemma-parametrize-solutions}
and Remark \ref{defos-remark-parametrize-solutions}
there exists an element
$$
\xi' \in
\Ext^1_A(\NL_{A/A'}, N) =
\Ext^1_B(\NL_{A/A'} \otimes_A^\mathbf{L} B, N) =
\Ext^1_B(L_{A/A'} \otimes_A^\mathbf{L} B, N)
$$
(for the equalities see Deformation Theory, Remark
\ref{defos-remark-parametrize-solutions} and use that
$\NL_{A'/A} = \tau_{\geq -1} L_{A'/A}$)
such that a solution exists if and only if this element is in
the image of the map
$$
\Ext^1_B(\NL_{B/A'}, N) = \Ext^1_B(L_{B/A'}, N)
\longrightarrow
\Ext^1_B(L_{A/A'} \otimes_A^\mathbf{L} B, N)
$$
The distinguished triangle (\ref{equation-triangle})
for $A' \to A \to B$ gives rise to a long exact sequence
$$
\ldots \to
\Ext^1_B(L_{B/A'}, N) \to
\Ext^1_B(L_{A/A'} \otimes_A^\mathbf{L} B, N) \to
\Ext^2_B(L_{B/A}, N) \to \ldots
$$
Hence taking $\xi$ the image of $\xi'$ works.
\end{proof}

\end{document}
